\scanpage{41}
sur $S$ modulo une des opérations $\psi_r$ ($r\in\mathbb Z_n^{*}$, ou $r\in\Psi_n$) ; \deleted{\unsure{Il revient donc au $=$|Je dis}} \inserted{en d'autres} termes, c'est une orbite $\pi$ de $\Psi_n$ opérant sur $\operatorname{Pol}(S)$. Regardant le ss-groupe \inserted{$Z$ de $\unsure{\mathfrak S_S|\mathcal S_S|G_S}$} engendré par un $\omega\in\pi$ (lequel ne dépend que de $\pi$, \deleted{\unsure{on est en|étant}} \inserted{étant} invariant en remplaçant $\omega$ par $\psi_r(\omega)$), on trouve un ss-groupe cyclique d'ordre $n$, et $S$ est un torseur sous ce groupe. Ainsi, \inserted{la donnée d'}une structure quasi-polygonale \deleted{sur} $S$ équivaut à la donnée d'une structure de torseur sur $S$ à groupe de \deleted{\unsure{gr|des}} translations \inserted{$Z$} cyclique (d'ordre $n$ nécessairement). Se donner un $\omega\in\pi$ revient au $=$ que de se donner un gén. ``au signe près'' dans $Z$. Se donner une structure polygonale \uline{orientée} compatible avec la structure quasi-polygonale revient à la donnée d'un gén. de $\unsure{Z|\mathbb Z}$.

La catégorie des \deleted{quasi} quasi-polygônes combinatoires d'ordre $n$ ($n\in\mathbb N$, $n\geq3$) est équivalente à celle des fibrés en droites affines sur $\operatorname{Spec}(\mathbb Z/n\mathbb Z)$, à un tel
% archive sheet number: 40
